\section{Backward Equations and Represented Formats}

\subsection{Exact attention backward}

For one batch/head, let
\begin{align}
  Z &= \alpha QK^\mathsf{T}, &
  P &= \softmax(Z+M), &
  O &= PV,
  \label{eq:exact-forward}
\end{align}
where $M_{ij}=0$ for legal entries and $-\infty$ for masked entries.  Given
the upstream gradient $G=\partial\mathcal L/\partial O$, exact backward is
\begin{align}
  dV &= P^\mathsf{T}G, \label{eq:exact-dv}\\
  dP &= GV^\mathsf{T}, \label{eq:exact-dp}\\
  r_i &= \sum_j P_{ij}dP_{ij}
       = \langle O_i,G_i\rangle, \label{eq:exact-row-dot}\\
  dS &= P\odot(dP-r\mathbf 1^\mathsf{T}), \label{eq:exact-ds}\\
  dQ &= \alpha dS K, \label{eq:exact-dq}\\
  dK &= \alpha dS^\mathsf{T}Q. \label{eq:exact-dk}
\end{align}
The identity in Equation~\ref{eq:exact-row-dot} lets the kernel compute the
row correction from saved O and $G$ before score reconstruction.  Causal
masking zeros all $j>i$ entries in P and therefore in $dS$.

As in FlashAttention, the implementation tiles the sequence dimensions and
recomputes scores and probabilities from Q, K, and the saved log-sum-exp
$L_i$; no $S\times S$ matrix is materialized in HBM
\cite{dao2022flashattention,zadouri2026flashattention4}.

\subsection{Quantizer notation}

Let $Q_F(x)$ be round-to-nearest, saturating conversion to format $F$.  For a
positive encode multiplier $m$, define
\begin{equation}
  \bar x_F = Q_F(mx),\qquad
  \widehat x_F=m^{-1}\bar x_F.
  \label{eq:quantizer-notation}
\end{equation}
An overbar denotes the stored code interpreted as its small floating value;
a hat denotes the decoded represented value.  E2M1 has the nonnegative
magnitudes
\begin{equation}
  \mathcal F_{\mathrm{E2M1}}^+
  =\left\{0,\tfrac12,1,\tfrac32,2,3,4,6\right\},
  \label{eq:e2m1-values}
\end{equation}
plus sign.  E4M3 is used both as a payload format and as the scale format for
NVFP4.  MXFP4 instead pairs E2M1 blocks of 32 values with E8M0 power-of-two
scales \cite{ocp2023mx,nvidia2025nvfp4}.

\subsection{Robust per-head FP4 Q/K}

For batch/head index $h$, define
\begin{align}
  a_h &=
    \max\!\left(\lVert Q_h\rVert_\infty,\lVert K_h\rVert_\infty\right),\\
  \rho_h &=
    \max\!\left(
      \sqrt{\frac{\lVert Q_h\rVert_F^2}{S d}},
      \sqrt{\frac{\lVert K_h\rVert_F^2}{S d}}
    \right).
  \label{eq:head-stats}
\end{align}
The retained clipping magnitude is
\begin{equation}
  c_h=\max(0.325\,a_h,\;2.75\,\rho_h).
  \label{eq:robust-clip}
\end{equation}
With implementation bounds $m_{\min}$ and $m_{\max}$, the target multiplier
is
\begin{equation}
  m_h^\star=
  \begin{cases}
    \operatorname{clip}\!\left(6/c_h,m_{\min},m_{\max}\right),
      &c_h>0,\\
    m_{\max},&c_h=0.
  \end{cases}
  \label{eq:target-multiplier}
\end{equation}
The tensor instruction consumes an E4M3 dequantizer.  The producer therefore
snaps the reciprocal first,
\begin{equation}
  \delta_h=Q_{\mathrm{E4M3}}\!\left((m_h^\star)^{-1}\right),
  \qquad m_h=\delta_h^{-1},
  \label{eq:snap-dequantizer}
\end{equation}
then packs
\begin{align}
  \bar Q_h &= Q_{\mathrm{E2M1}}(m_hQ_h),&
  \bar K_h &= Q_{\mathrm{E2M1}}(m_hK_h),\\
  \widehat Q_h &= \delta_h\bar Q_h,&
  \widehat K_h &= \delta_h\bar K_h.
  \label{eq:adaptive-qk}
\end{align}
Q and K deliberately share $m_h$.  Separate multipliers would require more
score scale state and a different instruction-facing layout.

\subsection{Full-rank FP4 score reconstruction}

$D_{QK}=192$ is split into three K64 chunks.  The represented base-two score
is
\begin{equation}
  \widetilde Z_2 =
  \alpha\log_2(e)
  \sum_{c=0}^{2}
  \mma_{\mathrm{NV4}}\!\left(
    \bar Q_{h,c},\bar K_{h,c}^{\mathsf T};
    \delta_h,\delta_h
  \right),
  \label{eq:represented-score}
\end{equation}
where the semicolon denotes the E4M3 scale page consumed by the block-scaled
instruction.  All three chunks are retained.  The two-command D128
approximation is disabled.

Let $L_i$ be the log-sum-exp saved by the forward that produced O.  The
backward probability used by the represented operator is
\begin{equation}
  \widetilde P_{ij}
  =2^{\widetilde Z_{2,ij}-L_i\log_2(e)}\,\mathbf 1_{\{j\le i\}}.
  \label{eq:represented-p}
\end{equation}
Because $L_i$ belongs to the forward operator while the score is reconstructed
from represented Q/K, Equation~\ref{eq:represented-p} is not, in general,
the exact softmax of $\widehat Q\widehat K^\mathsf T$.  This is the ordinary
saved-LSE backward contract, now evaluated with low-precision Q/K.

\subsection{Fixed-power E4M3 path}

The retained producer uses exact powers of two so scale work becomes exponent
adjustment.  Define
\begin{align}
  \bar G_4 &= Q_{\mathrm{E4M3}}(2^2G),&
  \bar V_4 &= Q_{\mathrm{E4M3}}(2^2V),&
  \bar P_8 &= Q_{\mathrm{E4M3}}(2^8\widetilde P).
  \label{eq:fp8-operands}
\end{align}
The subscripts 4 and 8 indicate the power-of-two multiplier, not bit width.
The $dP$ tensor instruction produces
\begin{equation}
  D_P=\mma_{\mathrm{E4M3}}(\bar G_4,\bar V_4^\mathsf T)
  \approx 2^4\,dP.
  \label{eq:scaled-dp}
\end{equation}
The mandatory statistics pass publishes the matching row term directly:
\begin{equation}
  R_i=2^4\langle O_i,G_i\rangle.
  \label{eq:scaled-row-term}
\end{equation}
Thus $D_P-R\mathbf1^\mathsf T$ can be centered without first dequantizing the
whole $dP$ tile.

For $dV$, the same $2^2$ E4M3 $G$ operand is reused:
\begin{equation}
  \widetilde{dV}
  =2^{-10}\,
    \mma_{\mathrm{E4M3}}(\bar P_8^\mathsf T,\bar G_4).
  \label{eq:represented-dv}
\end{equation}
The $2^{-10}$ correction is exactly $2^{-(8+2)}$.  This is why the retained
path does not need a second $2^8G$ representation.

\subsection{Fused E4M3 dS and mixed FP8-by-FP4 gradients}

E4M3 P is unpacked to FP16 pairs, while scaled $dP$ and $R$ are centered in
FP16.  Abstracting that local rounding as $\operatorname{fl}_{16}$, the
published dS code is
\begin{equation}
  \overline{dS}_{12}
  =Q_{\mathrm{E4M3}}\!\left(
    \operatorname{fl}_{16}(\bar P_8)\odot
    \operatorname{fl}_{16}(D_P-R\mathbf1^\mathsf T)
  \right)
  \approx Q_{\mathrm{E4M3}}(2^{12}dS).
  \label{eq:fused-ds}
\end{equation}
The factors $2^8$ and $2^4$ multiply to the required fixed dS scale
$s_{dS}=2^{12}=4096$.  Conversion is performed quarter by quarter from the
FP32 producer fragment into both instruction-facing orientations.  No
complete BF16 dS tile exists.

Finally,
\begin{align}
  \widetilde{dQ}
  &=\frac{\alpha}{2^{12}m_h}\,
    \mma_{\mathrm{E4M3}\times\mathrm{E2M1}}
    (\overline{dS}_{12},\bar K_h),
    \label{eq:represented-dq}\\
  \widetilde{dK}
  &=\frac{\alpha}{2^{12}m_h}\,
    \mma_{\mathrm{E4M3}\times\mathrm{E2M1}}
    (\overline{dS}_{12}^{\mathsf T},\bar Q_h).
    \label{eq:represented-dk}
\end{align}
The scalar corrections are fused into the existing dQ publication and dK
epilogue.  FP32 tensor accumulation is preserved before BF16 publication.

\subsection{Seven-word scale record}

Each batch/head has one FP32 record:
\begin{equation}
  \mathcal M_h=
  \left[
  m_h,\ m_h,\
  \frac{\alpha}{2^{12}m_h},\
  \frac{\alpha}{2^{12}m_h},\
  \frac{\alpha\log_2e}{m_h^2},\
  u_h,\
  \operatorname{bits}(\delta_h,\delta_h,\delta_h,\delta_h)
  \right].
  \label{eq:metadata-record}
\end{equation}
The first two words drive Q/K packing; words 2 and 3 are the live dQ/dK
corrections; word 5 is producer completion scratch; word 6 is the repeated
E4M3 scale-page byte pattern.  Word 4 keeps the explicit scalar score factor
in the interface contract, while the current hot score path applies the
equivalent $\delta_h$ scale page plus its ordinary $\alpha\log_2e$ factor.
